% from ega1-10-fr.tex
\oldpage[I]{201}

\begin{remark}[10.9.10]
\label{I.10.9.10-fr}
Si $X$ est la somme $X_1\amalg X_2$
(\hyperref[subsection:I.3.1-fr]{3.1}), $X'$ la réunion
$X'_1\cup X'_2$, où $X'_i$ est une partie fermée de $X_i$ ($i=1,2$),
on voit aussitôt que l'on a
$X_{/X'}={X_1}_{/X'_1}\amalg {X_2}_{/X'_2}$.
\end{remark}

\subsection{Application aux faisceaux cohérents sur les schémas formels
affines.}
\label{subsection:I.10.10-fr}

\begin{env}[10.10.1]
\label{I.10.10.1-fr}
Dans tout ce paragraphe, $A$ désignera un \emph{anneau adique
noethérien}, $\mathfrak{J}$ un idéal de définition de $A$. Soit
$X=\operatorname{Spec}(A)$, $\mathfrak{X}=\operatorname{Spf}(A)$, qui
s'identifie à la partie fermée $V(\mathfrak{J})$ de $X$
(\hyperref[I.10.1.2-fr]{10.1.2}). En outre, la définition
(\hyperref[I.10.1.2-fr]{10.1.2}) et la définition
(\hyperref[I.10.8.4-fr]{10.8.4}) montrent que le
\emph{schéma formel affine $\mathfrak{X}$} est identique au
\emph{complété $X_{/\mathfrak{X}}$} du schéma affine $X$ le long de la
partie fermée $\mathfrak{X}$ de son espace sous-jacent. A tout
$\mathscr{O}_X$-Module cohérent $\mathscr{F}$ correspond donc un
$\mathscr{O}_{\mathfrak{X}}$-Module de type fini
$\mathscr{F}_{/\mathfrak{X}}$, qui est d'ailleurs un faisceau de modules
topologiques sur le faisceau d'anneaux topologiques
$\mathscr{O}_{\mathfrak{X}}$. Mais tout $\mathscr{O}_X$-Module cohérent
$\mathscr{F}$ est de la forme $\widetilde M$, où $M$ est un $A$-module
de type fini (\hyperref[I.1.5.1-fr]{1.5.1})~; nous poserons
$(\widetilde M)_{/\mathfrak{X}}=M^\Delta$. En outre, si $u:M\to N$ est
un $A$-homomorphisme de $A$-modules de type fini, il lui correspond un
homomorphisme $\widetilde u:\widetilde M\to\widetilde N$, et par suite
aussi un homomorphisme continu
$\widetilde u_{/\mathfrak{X}}:(\widetilde M)_{/\mathfrak{X}}\to
(\widetilde N)_{/\mathfrak{X}}$, que nous noterons $u^\Delta$. Il est
immédiat que
$(v\circ u)^\Delta=v^\Delta\circ u^\Delta$~; on a ainsi défini un
\emph{foncteur additif covariant $M^\Delta$} de la catégorie des
$A$-modules de type fini dans celle des
$\mathscr{O}_{\mathfrak{X}}$-Modules de type fini. Lorsque $A$ est un
anneau \emph{discret}, on a $M^\Delta=\widetilde M$.
\end{env}

\begin{proposition}[10.10.2]
\label{I.10.10.2-fr}
\medskip\noindent
\begin{enumerate}
  \item[\rm{(i)}] $M^\Delta$ est un foncteur exact en $M$, et il existe
    un isomorphisme canonique fonctoriel de $A$-modules
    $\Gamma(\mathfrak{X},M^\Delta)\xrightarrow{\sim}M$.
  \item[\rm{(ii)}] Si $M$ et $N$ sont deux $A$-modules de type fini, il
    existe des isomorphismes canoniques fonctoriels
    \[
      (M\otimes_A N)^\Delta
      \xrightarrow{\sim}
      M^\Delta\otimes_{\mathscr{O}_{\mathfrak{X}}}N^\Delta
      \tag{10.10.2.1}\label{I.10.10.2.1-fr}
    \]
    \[
      (\operatorname{Hom}_A(M,N))^\Delta
      \xrightarrow{\sim}
      \mathscr{H}\!om_{\mathscr{O}_{\mathfrak{X}}}
      (M^\Delta,N^\Delta)
      \tag{10.10.2.2}\label{I.10.10.2.2-fr}
    \]
  \item[\rm{(iii)}] L'application $u\to u^\Delta$ est un isomorphisme
    fonctoriel
    \[
      \operatorname{Hom}_A(M,N)
      \xrightarrow{\sim}
      \operatorname{Hom}_{\mathscr{O}_{\mathfrak{X}}}
      (M^\Delta,N^\Delta)
      \tag{10.10.2.3}\label{I.10.10.2.3-fr}
    \]
\end{enumerate}
\end{proposition}

L'exactitude de $M^\Delta$ résulte de l'exactitude des foncteurs
$\widetilde M$ (\hyperref[I.1.3.5-fr]{1.3.5}) et
$\mathscr{F}_{/\mathfrak{X}}$ (\hyperref[I.10.8.8-fr]{10.8.8}). Par
définition, $\Gamma(\mathfrak{X},M^\Delta)$ est le séparé complété du
$A$-module $\Gamma(X,\widetilde M)=M$ pour la topologie
$\mathfrak{J}$-préadique~; mais comme $A$ est complet et $M$ de type
fini, on sait (\hyperref[0.7.3.6-fr]{0, 7.3.6}) que $M$ est séparé et
complet, ce qui achève de prouver (i). L'isomorphisme
(\hyperref[I.10.10.2.1-fr]{10.10.2.1}) (resp.
\hyperref[I.10.10.2.2-fr]{10.10.2.2}) provient de la composition des
isomorphismes (\hyperref[I.1.3.12-fr]{1.3.12, (i)}) et
(\hyperref[I.10.8.10.1-fr]{10.8.10.1}) (resp.
(\hyperref[I.1.3.12-fr]{1.3.12, (ii)}) et
(\hyperref[I.10.8.10.2-fr]{10.8.10.2})). Enfin, comme
$\operatorname{Hom}_A(M,N)$ est un $A$-module de type fini, on peut lui
appliquer (i), qui identifie
$\Gamma(\mathfrak{X},(\operatorname{Hom}_A(M,N))^\Delta)$ à
$\operatorname{Hom}_A(M,N)$, et utiliser
(\hyperref[I.10.10.2.2-fr]{10.10.2.2}), qui prouve que
l'homomorphisme (\hyperref[I.10.10.2.3-fr]{10.10.2.3}) est un
isomorphisme.

On déduit de (\hyperref[I.10.10.2-fr]{10.10.2}) toute une série de
conséquences analogues à celles déduites de
(\hyperref[I.1.3.7-fr]{1.3.7}) et
(\hyperref[I.1.3.12-fr]{1.3.12}), que nous laissons au lecteur le soin
de formuler.
